\SourceLeaf{382}%
\hspace*{\parindent}Je weniger sich hier der anomalische Gang verkennen lässt, desto interessanter wird die Frage, ob die astronomisch bestimmten Azimuthe der Dreiecksseiten ähnliche Anomalien zeigen. In der That steht, nach einem von \textsc{Laplace} zwar unter speciellen Beschränkungen aufgestellten, aber einer grossen Generalisirung fähigen Theorem, die Convergenz der Meridiane in einem nothwendigen und von der Gestalt der Erde unabhängigen Zusammenhange mit dem Längenunterschiede. so dass die Ungleichförmigkeiten der einen sich aus denen der andern, beim Fortschreiten in einer Kette von geodaetischen Linien, \textit{a priori} berechnen lassen. Da, wie wir berichtet haben, die astronomischen Azimuthalbestimmungen an den vier Hauptplätzen, Mailand, Turin, Mont Cenis und Colombier mit vieler Sorgfalt gemacht waren, so haben die Verf. mit diesen Orientirungen an den drei letzten Plätzen diejenigen verglichen, welche die Übertragung der Orientirung in Mailand vermittelst der geodaetischen Messungen ergibt, und dabei dieselben vorhin angezeigten Dimensionen des Erdsphäroids zum Grunde gelegt. Die Differenzen sind
\begin{center}\begin{tabular}{lr}
für Turin & \(-5''5\)\\
Mont Cenis & \(-51,2\)\\
Colombier & \(-25,2\)
\end{tabular}\end{center}
Auch hier erkennt man also ungemein grosse Anomalien. Allein wenn man nach dem erwähnten Theorem daraus die Anomalien der Längenunterschiede berechnet (was durch Division mit dem funfzehnfachen Sinus der Breite von Mailand und Veränderung des Zeichens geschieht), so ergeben sich Werthe, die von den unmittelbar gefundenen ganz verschieden sind, nemlich
\begin{center}\begin{tabular}{l|r|r}
& \multicolumn{1}{c|}{berechnete} & \multicolumn{1}{c}{Unterschied von der}\\
& \multicolumn{1}{c|}{Anomalie} & \multicolumn{1}{c}{beobacht. Anomalie}\\
Turin & \(+0''52\) & \(+2''60\)\\
Mont Cenis & \(+4,81\) & \(+4,11\)\\
Colombier & \(+2,34\) & \(+1,95\)
\end{tabular}\end{center}
Die Verf. bemerken über \textit{diese} Unterschiede bloss, dass sie zu gross seien, um der Anhäufung der Fehler bei den Winkelmessungen zur Last gelegt werden zu können, und lassen uns also im Dunkeln darüber, was wir von ihnen denken sollen. Nach unserer Ansicht sind diese drei Zahlen insofern von grösster Wich%
